% Class-compatibility preamble retained for a reader excerpt of
%   Dave Witte Morris, "A Cohomological Proof that Real Representations of
%   Semisimple Lie Algebras Have Q-Forms", SIGMA 11 (2015) 034,
%   DOI 10.3842/SIGMA.2015.034, CC BY 4.0.
%
% The source pins the publisher class sigma, which is not part of the delivery
% baseline, so this file supplies the amsart baseline and then copies the author's
% own declaration block mechanically and unchanged:
%   - pinned-original-source.tex lines 11-31: the author's DeclareMathOperator
%     list, the equation numbering and the author's newtheorem declarations;
%   - pinned-original-source.tex lines 35-43, which the author writes immediately
%     after \begin{document} in the publisher layout: the author's own \Res,
%     \ResOne and \Lie macros and \allowdisplaybreaks.
% The publisher header commands of the sigma class (\FirstPageHeading,
% \ArticleName and the like) and the macro \PaperNumber they require are not
% reproduced, because this excerpt carries neither the journal front matter nor
% anything else that uses them.
% No mathematical macro, statement, symbol or assumption is introduced, changed
% or removed by this file.  The author's full DeclareMathOperator list of pinned
% source lines 11-19 is carried, including GL, SL, SO and SU: these four are not
% supplied by the delivery baseline, so leaving them out makes the author's own
% notation (for instance GL_n(Q) in the corollary on lattices) undefined.
\documentclass[10pt,a4paper]{amsart}
\usepackage{amssymb,amsmath,amsthm}
\usepackage{enumerate}
\usepackage{hyperref}

\DeclareMathOperator{\Cliff}{Clif\/f}
\DeclareMathOperator{\Gal}{Gal}
\DeclareMathOperator{\GL}{GL}
\DeclareMathOperator{\Inn}{Inn}
\DeclareMathOperator{\Mat}{Mat}
\DeclareMathOperator{\Out}{Out}
\DeclareMathOperator{\SL}{SL}
\DeclareMathOperator{\SO}{SO}
\DeclareMathOperator{\SU}{\mathbf{SU}}
\numberwithin{equation}{section}

\newtheorem{prop}{Proposition}[section]
\newtheorem{cor}[prop]{Corollary}
\newtheorem{lem}[prop]{Lemma}
{\theoremstyle{definition}
\newtheorem{defn}[prop]{Definition}
\newtheorem{notation}[prop]{Notation}
\newtheorem{rem}[prop]{Remark}
\newtheorem{ques}[prop]{Question}
\newtheorem{eg}[prop]{Example}
\newtheorem{erratum}[prop]{Erratum}}

\newcommand{\Res}[1]{\mathop{\mathbf{R}_{#1}}}
\newcommand{\ResOne}[1]{\mathop{\mathbf{R}^{(1)}_{#1}}}
\newcommand{\Lie}[1]{\mathfrak{\lowercase{#1}}}
\allowdisplaybreaks
\usepackage{cleveref}
\crefname{originalremark}{original remark}{original remarks}
\Crefname{originalremark}{Original remark}{Original remarks}
\makeatletter
\newlabel{OnlyZStar}{{2.4}{source}{}{sourcepointer-1}{}}
\newlabel{OnlyZStar@cref}{{[originalremark][0][]2.4}{[1][0][]source}{}{}{}}
\makeatother
\title{R-universality of an odd-dimensional special orthogonal Lie algebra in terms of the even Clifford algebra}
\author{Dave Witte Morris}
\date{doi:10.3842/SIGMA.2015.034; CC BY 4.0}
\begin{document}
\maketitle
\noindent\textit{Editorial note.} All delivered mathematical text is the author's own, copied from the pinned publisher source sources/algebra-fl02-a-cohomological-proof-that-real-re/pinned-original-source.tex; no statement or proof step was rewritten, reordered or removed.\par\medskip
\noindent\textit{Editorial note.} Preamble provenance: the source pins the publisher class sigma, which is not part of the delivery baseline, so this card uses sources/algebra-fl02-a-cohomological-proof-that-real-re/retained-preamble.tex, which supplies the amsart baseline, copies the author's DeclareMathOperator list, equation numbering and newtheorem declarations from pinned source lines 11-31, and copies the author's own \textbackslash{}Res, \textbackslash{}ResOne and \textbackslash{}Lie macros from pinned source lines 35-43. No macro was invented and no mathematical symbol was changed.\par\medskip
\noindent\textit{Editorial note.} Bibliography provenance: the citations of the counted proof are resolved with sources/algebra-fl02-a-cohomological-proof-that-real-re/bibliography.tex, an excerpt of the author's own thebibliography environment (pinned source lines 954-1016) copied character for character, all entries kept so that every citation of a delivered excerpt resolves; no entry was rewritten, reordered, merged, split or invented and no citation key was changed.\par\medskip
\noindent\textit{Editorial note.} External references are navigation pointers only; every deliberately kept step of the author's argument stays inside the counted spans. Their printed numbers are the author's own numbering under the author's newtheorem declarations, computed mechanically by the lane's scanner.\par\medskip
\noindent\textit{Allowed standard prerequisites (editorial).}\begin{itemize}
\item The even Clifford algebra of a symmetric matrix over a field, and its Brauer class. Used inside the counted proof at source line 640 with the author's attribution to page 378 of the cited book KnusEtAl-BookOfInvols. Provides the Brauer class that the Tits algebra of the natural representation is equivalent to.
\item The remark that in the odd case the centre of the spin group has order two, so that only one Tits algebra has to be considered. The author's remark labelled OnlyZStar, invoked inside the counted proof at source line 639. Reduces the criterion to a single Brauer class, which is what makes the test one-dimensional.
\end{itemize}\medskip
\ifdefined\hypertarget\hypertarget{sourcepointer-1}{}\fi
\noindent\textit{Source reference.} Original remark 2.4: sources/algebra-fl02-a-cohomological-proof-that-real-re/pinned-original-source.tex, source lines 462--462. Navigation pointer only: the counted proof invokes this remark to reduce to a single Tits algebra; its content is declared as a named standard prerequisite.\par\medskip
\begin{prop}
Let $\Lie G = \Lie{SO}_n(B; \mathbb{Q})$, where~$B$~is a~symmetric matrix in $\GL_n(\mathbb{Q})$, and~$n$~is odd.
Then $\Lie G$ is $\mathbb{R}$-universal if and only if either $\Cliff^0_\mathbb{Q}(B)$ is split or
$\Cliff^0_\mathbb{R}(B)$ is not split.
\end{prop}

\begin{proof}
Since the center of $\mathbf{Spin}_n$ has order~$2$ when~$n$~is odd, there is only one Tits algebra to consider
(see Remark~\ref{OnlyZStar}), and it is Brauer equivalent to $\Cliff^0_F(B)$~\cite[p.~378]{KnusEtAl-BookOfInvols}.
\end{proof}

%% Bibliography excerpt retained for reader excerpts of
%   Dave Witte Morris, "A Cohomological Proof that Real Representations of Semisimple Lie Algebras Have Q-Forms",
%   SIGMA 11 (2015) 034, DOI 10.3842/SIGMA.2015.034.
%%
%% Every entry below is copied character for character from the author's own
% thebibliography environment of the pinned publisher source (lines 954-1016).
%% No entry was rewritten, reordered, merged, split or invented, and no citation
% key was changed.  All 12 entries are kept so that any citation of a delivered
%% excerpt resolves without guessing an entry.
\begin{thebibliography}{99}

\footnotesize\itemsep=0pt

\bibitem{DemazureGrothendieck-SGA3}
Demazure M., Grothendieck A., Sch\'emas en groups.~III. Structure des sch\'emas
  en groupes reductifs, S\'eminaire de g\'eom\'etrie alg\'ebrique du {B}ois
  {M}arie, 1962/1964, available at
  \url{http://library.msri.org/books/sga/sga/pdf/sga3-3.pdf}.

\bibitem{Harder-BerichtNeuere}
Harder G., Bericht \"uber neuere {R}esultate der {G}aloiskohomologie
  halbeinfacher {G}ruppen, \textit{Jber. Deutsch. Math.-Verein.} \textbf{70}
  (1968), 182--216, available at
  \url{http://www.digizeitschriften.de/en/dms/img/?PPN=PPN37721857X_0070&DMDID=dmdlog17}.

\bibitem{Kneser-LectGalCohoClassical}
Kneser M., Lectures on {G}alois cohomology of classical groups, \textit{Tata
  Institute of Fundamental Research Lectures on Mathematics}, Vol.~47, Tata
  Institute of Fundamental Research, Bombay, 1969, available at
  \url{http://www.math.tifr.res.in/~publ/ln/tifr47.pdf}.

\bibitem{KnusEtAl-BookOfInvols}
Knus M.A., Merkurjev A., Rost M., Tignol J.P., The book of involutions,
  \textit{American Mathematical Society Colloquium Publications}, Vol.~44,
  Amer. Math. Soc., Providence, RI, 1998.

\bibitem{Lam-IntroQuadForms}
Lam T.Y., Introduction to quadratic forms over f\/ields, \textit{Graduate Studies
  in Mathematics}, Vol.~67, Amer. Math. Soc., Providence, RI, 2005.

\bibitem{Morris-RRepsQForms}
Morris D., Real representations of semisimple {L}ie algebras have {${\mathbb
  Q}$}-forms, in Algebraic Groups and Arithmetic, Tata Institute of Fundamental
  Research, Mumbai, 2004, 469--490.

\bibitem{PlatonovRapinchukBook}
Platonov V., Rapinchuk A., Algebraic groups and number theory, \textit{Pure and
  Applied Mathematics}, Vol.~139, Academic Press, Inc., Boston, MA, 1994.

\bibitem{PrasadRapinchuk-Prescribed}
Prasad G., Rapinchuk A., On the existence of isotropic forms of semi-simple
  algebraic groups over number f\/ields with prescribed local behavior,
  \href{http://dx.doi.org/10.1016/j.aim.2006.01.001}{\textit{Adv. Math.}} \textbf{207} (2006), 646--660.

\bibitem{Raghunathan-ArithLattsSSGrps}
Raghunathan M.S., Arithmetic lattices in semisimple groups, \href{http://dx.doi.org/10.1007/BF02967980}{\textit{Proc.
  Indian Acad. Sci. Math. Sci.}} \textbf{91} (1982), 133--138.

\bibitem{Tits-Classification}
Tits J., Classif\/ication of algebraic semisimple groups, in Algebraic {G}roups
  and {D}iscontinuous {S}ubgroups ({P}roc. {S}ympos. {P}ure {M}ath., {B}oulder,
  {C}olo., 1965), Amer. Math. Soc., Providence, R.I., 1966, 33--62.

\bibitem{Tits-RepLinIrred}
Tits J., Repr\'esentations lin\'eaires irr\'eductibles d'un groupe r\'eductif
  sur un corps quelconque, \href{http://dx.doi.org/10.1515/crll.1971.247.196}{\textit{J.~Reine Angew. Math.}} \textbf{247} (1971),
  196--220.

\bibitem{Tits-StronglyInner}
Tits J., Strongly inner anisotropic forms of simple algebraic groups,
  \href{http://dx.doi.org/10.1016/0021-8693(90)90201-X}{\textit{J.~Algebra}} \textbf{131} (1990), 648--677.


\end{thebibliography}
\end{document}
